\subsection{INITIAL UNIFORMITIES}

PROPOSITION 4. Let X be a set, let \((\mathrm{Y}_{t})_{t\in I}\) be a family of uniform spaces, and for each \(t\in I\) let \(f_{t}\) be a mapping of X into \(Y_{t}\) . For each \(t\in I\) let \(g_{t}\) denote \(f_{t}\times f_{t}\) . Let S be the set of subsets of \(X\times X\) of the form \(\overline{g}_{t}(V_{t})\) , where \(t\in I\) and \(V_{t}\) is an entourage of \(Y_{t}\) , and let B be the set of all finite intersections

\[
\mathbf {U} \left(\mathrm{V} _ {\iota_ {1}}, \dots , \mathrm{V} _ {\iota_ {n}}\right) = _ {i} ^ {- 1} g _ {i} \left(\mathrm{V} _ {\iota_ {1}}\right) \cap \dots \cap^ {- 1} g _ {\iota_ {n}} \left(\mathrm{V} _ {\iota_ {n}}\right)\tag{1}
\]

of sets of \(\mathfrak{S}\) . Then \(\mathfrak{B}\) is a fundamental system of entourages of a uniformity \(\mathfrak{U}\) on \(\mathbf{X}\) which is the initial uniform structure on \(\mathbf{X}\) with respect to the family \((f_{\mathfrak{i}})\) (Set Theory, Chapter IV, § 2, no. 3), and in particular \(\mathfrak{U}\) is the coarsest uniformity on \(\mathbf{X}\) for which all the mappings \(f_{\mathfrak{i}}\) are uniformly continuous. But otherwise, let \(h\) be a mapping of a uniform space \(\mathbf{Z}\) into \(\mathbf{X}\) ; then \(h\) is uniformly continuous (when \(\mathbf{X}\) is endowed with the uniformity \(\mathfrak{U}\) ) if and only if each of the mappings \(f_{\mathfrak{i}} \circ h\) is uniformly continuous.

It is immediately seen that \(\mathfrak{B}\) satisfies axioms \((\mathbf{B}_1)\mathbf{Y}\) and \((\mathbf{U}_1')\) . If \(\mathbf{W}_t = \overline{g}_t^1 (\mathbf{V}_t)\) , then \(\overline{\mathbf{W}}_t = \overline{g}_t^1 (\overline{\mathbf{V}}_t^1)\) and \(\dot{\mathbf{W}}_t = \overline{g}_t^1 (\overline{\mathbf{V}}_t^2)\) ; hence \(\mathfrak{B}\) also satisfies axioms \((\mathbf{U}_{\mathrm{II}}^t)\) and \((\mathbf{U}_{\mathrm{III}}^t)\) and is therefore a fundamental system of entourages of a uniformity \(\mathcal{U}\) on \(\mathbf{X}\) . Furthermore, it follows immediately from the definition of \(\mathcal{U}\) and Definition 1 and no. 1 that \(f_{t}\) is uniformly continuous for each index \(i\in I\) ; hence (no. 1, Proposition 2) \(f_{t}\circ h\) is uniformly continuous for each \(i\in I\) if \(h\) is. Conversely, suppose that \(f_{t}\circ h\) is uniformly continuous for each \(i\in I\) , and consider a set \(\mathrm{U}(\mathrm{V}_{i_t},\dots,\mathrm{V}_{i_n})\) ; by hypothesis, for each \(k\) such that \(1\leqslant k\leqslant n\) , there is an entourage \(\mathbf{W}_k\) of \(\mathbf{Z}\) such that the relation \((z,z^{\prime})\in \mathbf{W}_k\) implies \([f_{i_k}(h(z)),f_{i_k}(h(z')))]\in \mathbf{V}_k\) ; if

\[
\mathbf {W} = \bigcap_ {k} \mathbf {W} _ {k},
\]

these \(n\) relations are simultaneously satisfied whenever \(z\) and \(z'\) are W-close, so that we have then \((h(z), h(z')) \in \mathrm{U}(\mathrm{V}_{\iota_1}, \ldots, \mathrm{V}_{\iota_n})\) , and the proof is complete.

COROLLARY. The topology on X induced by the coarsest uniformity U for which the \(f_{i}\) are uniformly continuous is also the coarsest topology for which the \(f_{i}\) are continuous.

This is an immediate consequence of the definition of the neighbourhoods of a point in this latter topology (Chapter I, § 2, no. 3, Proposition 4).

The general properties of initial structures (Set Theory, Chapter IV, § 2, no. 3, criterion CST 10) imply in particular the following transitivity property:

PROPOSITION 5. Let X be a set, let \((Z_{\iota})_{\iota \in I}\) be a family of uniform spaces, let \((J_{\lambda})_{\lambda \in L}\) be a partition of I and let \((Y_{\lambda})_{\lambda \in L}\) be a family of sets indexed by L. For each \(\lambda \in L\) , let \(h_{\lambda}\) be a mapping of X into \(Y_{\lambda}\) ; for each \(\lambda \in L\) and each \(\iota \in J_{\lambda}\) , let \(g_{\iota \lambda}\) be a mapping of \(Y_{\lambda}\) into \(Z_{\iota}\) ; and let \(f_{\iota} = g_{\iota \lambda} \circ h_{\lambda}\) . Let each \(Y_{\lambda}\) carry the coarsest uniform structure for which the mappings \(g_{\iota \lambda} (\iota \in J_{\lambda})\) are uniformly continuous. Then the coarsest uniform structure on X for which the \(f_{\iota}\) are uniformly continuous is the same as the coarsest uniform structure on X for which the \(h_{\lambda}\) are uniformly continuous.

